\chapter{Evaluación experimental}
\label{cap:9}

Los capítulos anteriores han definido el problema, el marco normativo, la arquitectura
del sistema, el conjunto de datos y el procedimiento de construcción de etiquetas académicas
P1--P4. Este capítulo cierra el ciclo experimental evaluando la Capa 2 de priorización
interpretable sobre un flujo reproducible: conjunto limpio, etiquetas débiles,
particiones, línea base heurística, modelos RuleFit y controles contra fugas de
información.

La evaluación no pretende demostrar validez operativa en producción. El Registro 112
CyL no contiene una prioridad oficial histórica comparable con la recomendación del
sistema. Por tanto, la evaluación mide el comportamiento frente a la etiqueta académica
P1--P4 generada por supervisión débil y documenta sus limitaciones. Esta distinción es
central para no sobredimensionar los resultados.

\section{Objetivos de evaluación}

La evaluación experimental persigue cinco objetivos:

\begin{enumerate}
    \item Verificar que la Capa 2 aprende a reproducir la etiqueta académica P1--P4 con
    mejor rendimiento que una línea base heurística;
    \item Comprobar que la sensibilidad de P1 se mantiene por encima del criterio
    fijado en 0.85;
    \item Preservar la interpretabilidad mediante un límite de 30 reglas activas
    exportadas;
    \item Contrastar RuleFit-lite con una implementación externa de la misma familia en una prueba diagnóstica;
    \item Verificar que no se utilizan columnas posteriores a la decisión.
\end{enumerate}

Las métricas principales son exactitud, macro-F1, precisión, sensibilidad y F1 por clase,
sensibilidad P1, número de reglas activas, tiempo de entrenamiento e inferencia media por
fila. La macro-F1 se emplea como métrica de selección porque evita que el rendimiento
quede dominado por las clases más frecuentes. La sensibilidad P1 se prioriza por el
coste potencial de los falsos negativos en la clase crítica.

\section{Niveles de evaluación}

La evaluación se estructura en tres niveles complementarios. Primero, una evaluación
de Capa 2 sobre las particiones definidas, donde se comparan la línea base heurística,
RuleFit-lite y una implementación diagnóstica. Segundo, una prueba técnica integrada reproducible de
escenarios Capa 1--Capa 2--Capa 3, orientada a comprobar que la cadena completa conserva
coherencia funcional. Tercero, una prueba funcional integrada del prototipo con LLM local,
cuyo objetivo no es medir rendimiento estadístico, sino comprobar la ejecución
extremo a extremo en casos P1--P4.

\begin{figure}[htbp]
\centering
\small
\fbox{%
\begin{tabular}{c}
\textbf{Conjunto de Datos Limpio} (9.380 registros del Registro 112 CyL) \\
\(\downarrow\) \\
\textbf{Supervisión Débil} (4 funciones heurísticas correlacionadas \(\rightarrow\) Krippendorff \(\alpha = 0{,}899\)) \\
\(\downarrow\) \\
\textbf{Etiqueta Académica Proxy} (Clases P1--P4) \\
\(\downarrow\) \\
\textbf{Particiones Estratificadas} \\
(entrenamiento, validación, prueba y robustez temporal) \\
\(\downarrow\) \\
\textbf{Comparativa de Modelos} (línea base, RuleFit-lite e \texttt{imodels}) \\
\(\downarrow\) \\
\textbf{Evaluación Cuantitativa} \\
(métricas por clase, sensibilidad P1 y 59 falsos negativos) \\
\(\downarrow\) \\
\textbf{Fidelidad Explicativa} (seis criterios y 119 casos) \\
\(\downarrow\) \\
\textbf{Evaluación Exploratoria con Usuarios} \\
(nueve participantes; estudio descriptivo no probabilístico)
\end{tabular}
}
\caption{Secuencia metodológica del flujo de evaluación experimental.}
\label{fig:flujo-evaluacion-cap9}
\end{figure}

\section{Datos, etiquetas y particiones}

El experimento usa el conjunto limpio, formado por 9.380 registros del Registro 112 Castilla y León. La etiqueta objetivo P1--P4 se
obtiene mediante supervisión débil con cuatro funciones heurísticas correlacionadas que comparten una puntuación programática de gravedad. No se trata de anotadores independientes ni de modelos NER, agrupamiento semántico o LLM ejecutados durante el etiquetado. La ejecución principal
alcanza Krippendorff \(\alpha = 0{,}899902\), superior al umbral de \(0{,}67\). La
ablación sin reglas heurísticas mantiene \(\alpha = 0{,}834155\). Este resultado
demuestra que el acuerdo no desaparece al retirar una fuente, pero no garantiza
independencia entre las etiquetas y las variables predictivas.

Las particiones se generan de forma estratificada por año, provincia inferida y etiqueta
P1--P4, con semilla fija 42. La Tabla~\ref{tab:9-splits} resume los tamaños y la
distribución de etiquetas.

\begin{table}[ht]
\centering
\caption{Distribución de etiquetas por partición}
\label{tab:9-splits}
\small
\begin{tabular}{|l|r|r|r|r|r|}
\hline
\textbf{Partición} & \textbf{Registros} & \textbf{P1} & \textbf{P2} & \textbf{P3} & \textbf{P4} \\ \hline
Entrenamiento & 6.512 & 3.049 & 1.504 & 1.701 & 258 \\ \hline
Validación & 1.415 & 658 & 329 & 374 & 54 \\ \hline
Prueba & 1.453 & 655 & 349 & 375 & 74 \\ \hline
\end{tabular}
\end{table}

Adicionalmente se ha generado una partición temporal alternativa: entrenamiento hasta 2020,
validación en 2021 y prueba en 2022. Esta partición se utiliza como prueba de robustez
temporal, mientras que el experimento principal usa la partición estratificada por su mayor
equilibrio entre años, provincias y etiquetas.

\section{Política anti-leakage}

La evaluación excluye todas las columnas conocidas después de la decisión operativa o
derivadas de procesos post-hoc. En particular, no se usan campos relativos a medios movilizados, pacientes atendidos, cierre del incidente, actualizaciones posteriores, enlaces externos ni columnas auxiliares sin significado operativo.

Las variables permitidas se limitan a señales textuales previas a la decisión y a
variables categóricas derivadas de la categoría preliminar. La relación completa
se documenta en el Capítulo~\ref{cap:7}. Además, el control reproducible de fugas verifica que el motor de priorización no dependa de columnas prohibidas.

\section{Modelos evaluados}

Se comparan tres aproximaciones:

\begin{description}
    \item[Línea base heurística.] Conjunto de 19 reglas manuales con anclaje normativo,
    descritas en el Anexo~\ref{cap:anexoe}. Su función es actuar como línea base
    interpretable a batir.
    \item[RuleFit-lite.] Modelo seleccionado para Capa 2. Combina reglas de
    árboles poco profundos con una capa logística dispersa. Se ajusta con un submuestreo
    estratificado de 3.000 casos de los 6.512 disponibles y exporta 30 reglas activas.
    \item[RuleFit externo.] Implementación canónica de la familia RuleFit. Como la configuración utilizada trabaja en clasificación binaria,
    se emplea un esquema one-vs-rest P1--P4. Se ajusta con 50 casos y su papel es exclusivamente diagnóstico; no constituye una comparación equilibrada con RuleFit-lite.
\end{description}

\section{Selección del modelo}

La selección se realiza sobre la partición de validación de la configuración
experimental disponible. La Tabla~\ref{tab:9-comparativa} muestra que RuleFit-lite obtiene la
mejor macro-F1 y mantiene la sensibilidad P1 por encima del umbral de \(0{,}85\). No obstante,
los resultados de \texttt{imodels} deben leerse solo como referencia diagnóstica, dado que su
muestra de ajuste es mucho menor.

\begin{table}[ht]
\centering
\caption{Comparativa de modelos de Capa 2 en validación}
\label{tab:9-comparativa}
\small
\begin{tabular}{|l|r|r|r|}
\hline
\textbf{Modelo} & \textbf{Exactitud} & \textbf{Macro-F1} & \textbf{Sensibilidad P1} \\ \hline
Línea base heurística & \(0{,}717314\) & \(0{,}498785\) & \(0{,}995441\) \\ \hline
RuleFit-lite & \(0{,}881188\) & \(0{,}879592\) & \(0{,}934400\) \\ \hline
RuleFit externo & \(0{,}789250\) & \(0{,}584768\) & \(0{,}910400\) \\ \hline
\end{tabular}
\end{table}

La línea base obtiene una sensibilidad P1 muy elevada por su diseño conservador, pero discrimina
peor entre las cuatro clases. RuleFit-lite mejora la macro-F1 de validación y ofrece
30 reglas exportables. La alternativa externa presenta menor rendimiento en la ejecución
reducida disponible, pero la diferencia de tamaños de ajuste impide inferir superioridad
computacional o estadística de RuleFit-lite en condiciones equivalentes.

Por ello, RuleFit-lite se selecciona antes de consultar la evaluación final. La
partición de prueba canónica, regenerada posteriormente y formada por 1.453 casos, se
utiliza una sola vez para estimar el rendimiento del modelo elegido: exactitud
\(0{,}903648\), macro-F1 \(0{,}905604\) y sensibilidad P1 \(0{,}909924\).

Estas métricas corresponden al estimador RuleFit-lite evaluado sin el posprocesamiento
de la Puerta de Seguridad. La salvaguarda se aplica durante la inferencia integrada y se
analiza mediante escenarios funcionales, pero no se mezcla con la comparación estadística
entre modelos.

\section{Métricas por clase del modelo seleccionado}

La Tabla~\ref{tab:9-rulefit-lite-clases} resume las métricas por clase de RuleFit-lite
en la partición de prueba.

\begin{table}[ht]
\centering
\caption{Métricas por clase de RuleFit-lite en la partición de prueba}
\label{tab:9-rulefit-lite-clases}
\small
\begin{tabular}{|l|r|r|r|r|}
\hline
\textbf{Clase} & \textbf{Precisión} & \textbf{Sensibilidad} & \textbf{F1} & \textbf{Soporte} \\ \hline
P1 & 0.941548 & 0.909924 & 0.925466 & 655 \\ \hline
P2 & 0.822888 & 0.865330 & 0.843575 & 349 \\ \hline
P3 & 0.924119 & 0.909333 & 0.916667 & 375 \\ \hline
P4 & 0.880952 & 1.000000 & 0.936709 & 74 \\ \hline
\end{tabular}
\end{table}

El comportamiento por clase muestra un equilibrio razonable. P1 mantiene sensibilidad
y precisión elevadas, un aspecto relevante por el coste potencial de infrapriorizar esta clase. P2 es la
clase más difícil, con F1 inferior al resto, un resultado esperable al tratarse de una
zona frontera entre gravedad crítica y seguimiento ordinario. P4 obtiene sensibilidad total
en la partición de prueba, aunque con menor precisión por su baja frecuencia relativa.

\section{Matriz de confusión del modelo seleccionado}

La Tabla~\ref{tab:9-confusion-rulefit-lite} muestra la matriz de confusión de
RuleFit-lite en la partición de prueba.

\begin{table}[ht]
\centering
\caption{Matriz de confusión de RuleFit-lite en la partición de prueba}
\label{tab:9-confusion-rulefit-lite}
\small
\begin{tabular}{|l|r|r|r|r|}
\hline
\textbf{Referencia / Predicción} & \textbf{P1} & \textbf{P2} & \textbf{P3} & \textbf{P4} \\ \hline
P1 & 596 & 34 & 18 & 7 \\ \hline
P2 & 37 & 302 & 10 & 0 \\ \hline
P3 & 0 & 31 & 341 & 3 \\ \hline
P4 & 0 & 0 & 0 & 74 \\ \hline
\end{tabular}
\end{table}

Los 7 casos P1 clasificados como P4 constituyen el principal riesgo detectado en la
evaluación. Aunque el rendimiento global del modelo es alto, estos errores representan
infrapriorizaciones críticas y deben tratarse como prioridad de revisión humana,
análisis de causas y mejora futura del sistema. El resto de falsos negativos P1 se
distribuye en 34 casos P1--P2 y 18 casos P1--P3. La representación gráfica de esta matriz se conserva como evidencia reproducible de evaluación.

\section{Robustez temporal y sesgo interno}

Para comprobar que el rendimiento no depende exclusivamente de la partición estratificada, se
evalúa el mismo modelo sobre la partición temporal de 2022. En este escenario se obtienen
735 registros, exactitud 0.884354, macro-F1 0.877594 y sensibilidad P1 0.928767. La estabilidad
entre ambos cortes no muestra una degradación temporal relevante en los datos evaluados,
aunque no sustituye a una validación prospectiva con datos nuevos.

El análisis por subgrupos se realiza por provincia, año y categoría operativa. En
el reporte canónico, las provincias con menor macro-F1 son Segovia
\(0{,}887379\), León \(0{,}893354\) y Salamanca \(0{,}894117\). Por
categoría, el grupo sanitario presenta macro-F1 \(0{,}792966\) y sensibilidad P1
\(0{,}404255\); incendio obtiene \(0{,}861823\) y \(0{,}731707\); y tráfico,
\(0{,}935353\) y \(0{,}964727\), respectivamente. Estos resultados no demuestran
sesgo estructural, pero identifican los casos sanitarios y de incendio como focos
prioritarios de revisión. Las gráficas de apoyo se conservan como evidencia complementaria, pero la interpretación principal se recoge en este capítulo.

\section{Interpretabilidad}

RuleFit-lite exporta 30 reglas activas, cumpliendo el requisito de parsimonia. Las
reglas combinan señales legibles, por ejemplo riesgo vital textual, accidente de
tráfico, rescate, incendio o ausencia de señales críticas. Esta salida facilita
su uso por la Capa 3, que puede transformar las reglas activadas en una explicación
operativa.

La línea base heurística queda documentada en el Anexo~\ref{cap:anexoe}; la ficha de modelo de
Capa 2 y la comparativa completa se recogen en el Anexo~\ref{cap:anexol}. La evidencia reproducible de cierre enlaza métricas finales, partición temporal, tablas de análisis por subgrupos, falsos negativos P1, muestra de revisión y fidelidad explicativa.

\section{Revisión interna y trazabilidad}

La evaluación automática se complementa con una revisión interna conjunta realizada
por los tres autores. Su objetivo es examinar cada recomendación con evidencias suficientes,
centrando la atención en los errores críticos y en la trazabilidad de la decisión. Esta
revisión cualitativa no equivale a tres anotaciones independientes ni permite calcular
acuerdo interrevisor.

Cada ficha preparada muestra el texto operativo minimizado, provincia, fecha,
categoría preliminar, etiqueta académica, acuerdo de supervisión débil,
recomendación P1--P4, confianza, probabilidades por clase, reglas activadas,
variables permitidas, columnas excluidas por anti-\textit{leakage} y, cuando la Capa 3 está
disponible, la explicación normativa asociada. Esta ficha permite que el revisor acepte
la recomendación, la modifique o la rechace.

La plantilla permite registrar una futura decisión humana independiente asociada a la inferencia,
incluyendo identificador del caso, prioridad recomendada, prioridad final asignada,
motivo de divergencia, revisor, marca temporal y huella de entrada. En la revisión conjunta de los autores se utilizaron las fichas como instrumento cualitativo, sin emplear esos campos para construir una nueva referencia supervisada. Las divergencias de dos o más niveles se consideran casos de auditoría especial.

La muestra revisada incluye todos los falsos negativos P1 de la
partición de prueba estratificada, junto con casos frontera P1/P2, P1/P3, P1/P4 y P2/P3. En la
ejecución final se detectan 59 falsos negativos P1: 34 P1--P2, 18 P1--P3 y 7 P1--P4.
La muestra contiene 119 casos: los 59 falsos negativos P1 y 60 casos equilibrados, 15 por cada clase P1--P4. La revisión conjunta aporta una auditoría interna centrada en los errores críticos, pero no constituye validación externa con operadores del 112 ni una evaluación independiente entre revisores.

\section{Fidelidad de explicaciones}

Sobre la misma muestra de 119 casos se ejecuta una evaluación automatizada y local
de fidelidad. El objetivo es comprobar que la Capa 3 explica
lo que realmente entrega la Capa 2: prioridad recomendada, reglas o señales activadas,
probabilidades y necesidad de cautela. La evaluación se realiza mediante un juez determinista local y se conserva como evidencia reproducible de fidelidad.

La versión ejecutada utiliza un juez determinista reproducible, no un LLM-as-Judge
externo. La Capa 3 se evalúa como sistema de fidelidad explicativa: la explicación
generada debe ser coherente con la prioridad y reglas producidas por Capa 2. No se
evalúa Capa 3 como modelo alternativo de decisión. La rúbrica verifica seis criterios:
alineación con la prioridad recomendada, trazabilidad de reglas, cobertura de señales,
citas legales en P1/P2, ausencia de contradicciones y presencia de disclaimer de
confianza o respuesta alternativa determinista. Sobre los 119 casos se obtiene una tasa de cumplimiento de
\(1{,}000000\), fidelidad media \(0{,}957367\) y mínima \(0{,}933333\). En los
59 falsos negativos P1, la fidelidad media es \(0{,}952494\).

Estos resultados no prueban calidad lingüística ni aceptación por operadores, pero sí
aportan evidencia de que, en la muestra evaluada y con respuesta alternativa determinista,
las explicaciones no contradicen la salida de Capa 2 y conservan la trazabilidad definida. Una
evaluación futura con LLM-as-Judge independiente, en la línea propuesta por
\citeA{Zheng2023LLMJudge}, puede reutilizar los mismos 119 casos.

\section{Posicionamiento frente a trabajos relacionados}

La Tabla~\ref{tab:9-comparativa-autores} contrasta el enfoque propuesto con los trabajos relacionados del estado del arte, siguiendo la recomendación de presentar resultados propios frente a otros autores.

\begin{table}[htbp]
\centering
\caption{Comparativa del enfoque propuesto frente a trabajos relacionados}
\label{tab:9-comparativa-autores}
\scriptsize
\setlength{\tabcolsep}{3pt}
\begin{tabularx}{\textwidth}{|p{2.5cm}|p{1.7cm}|p{1.7cm}|p{1.4cm}|p{1.3cm}|p{1.4cm}|X|}
\hline
\textbf{Enfoque} & \textbf{Datos reales} & \textbf{Interp.} & \textbf{Expl. norm.} & \textbf{On-prem.} & \textbf{Superv.} & \textbf{Métrica clave} \\ \hline
AIDR~\cite{Imran2014AIDR} & Redes sociales & No (DNN) & No & No & Fuerte & Acc. multi-label \\ \hline
CrisisTransf. \cite{Lamsal2024} & Crisis real & No (transformer) & No & Parcial & Fuerte & F1 clasificación \\ \hline
TREC-IS\\~\cite{TRECIS2021} & Twitter & No (ensemble) & No & No & Fuerte & Benchmarks \\ \hline
Línea base heurística (este TFM) & 112 CyL & Sí (reglas) & Parcial & Sí & Débil & F1 0{,}464; sensibilidad P1 0{,}998 \\ \hline
\textbf{RuleFit-lite + Capa 3} & \textbf{112 CyL} & \textbf{Sí} & \textbf{Sí} & \textbf{Sí} & \textbf{Débil} & \textbf{F1 0{,}906; P1 0{,}910} \\ \hline
\end{tabularx}
\smallskip
{\small\textit{Fuente: elaboración propia a partir de los trabajos citados.}\par}
\end{table}

\smallskip
Las ventajas diferenciales del enfoque son: (1)~datos empíricos reales del 112 CyL frente a redes sociales; (2)~núcleo interpretable con reglas auditables; (3)~explicación normativa local sin APIs externas; y (4)~supervisión débil que evita la dependencia de etiquetado manual masivo. La limitación principal es la etiqueta académica P1--P4, no validada por personal operativo.

\section{Pruebas integradas del prototipo}

Además de la evaluación offline de Capa 2, se han ejecutado dos pruebas integradas del
prototipo. La primera comprueba la cadena reproducible Capa 1--Capa 2--Capa 3 con
respuesta alternativa determinista sobre escenarios curados. La segunda verifica la ejecución completa con
LLM local. Ninguna sustituye a las métricas de la partición de prueba ni a la validación externa, pero
ambas verifican que la cadena software real funciona de extremo a extremo: texto del
incidente, extracción de señales en Capa 1, priorización RuleFit-lite en Capa 2,
explicación en Capa 3 y respuesta consumible por la interfaz.

Tras integrar la Capa 1, el sistema supera 5/5 escenarios curados con RuleFit-lite activo:
incendio urbano con atrapados, accidente de tráfico con inconsciente, fuga química con
intoxicación, inundación sin víctimas e incidencia vial ordinaria. Este último caso
incluye negaciones explícitas (``no hay heridos ni vehículos atrapados'') y confirma
que Capa 1 no activa falsos positivos de atrapamiento tras el saneamiento realizado.

La ejecución confirma la disponibilidad del LLM local. Los cuatro casos representativos P1--P4 se
clasifican conforme a la prioridad esperada, con RuleFit activo y sin recurrir al modo
alternativo determinista de Capa 3. La latencia media observada es \(4{,}90\) s por inferencia
integrada.

\begin{table}[ht]
\centering
\caption{Resultado de la prueba funcional integrada P1--P4 con LLM local}
\label{tab:9-smoke-prototipo}
\small
\begin{tabularx}{\textwidth}{|p{3.6cm}|p{1.7cm}|p{1.7cm}|X|}
\hline
\textbf{Caso} & \textbf{Esperada} & \textbf{Obtenida} & \textbf{Distribución principal} \\ \hline
Incendio con atrapados & P1 & P1 & P1 97.35\%, P2 2.61\% \\ \hline
Fuga de gas sin heridos & P2 & P2 & P2 92.21\%, P3 5.03\% \\ \hline
Árbol en calzada & P3 & P3 & P3 93.71\%, P2 3.79\% \\ \hline
Consulta vecinal sin riesgo & P4 & P4 & P4 98.54\%, P3 1.40\% \\ \hline
\end{tabularx}
\end{table}

Esta evidencia permite separar la evaluación cuantitativa de Capa 2, ya realizada sobre la partición de prueba estratificada y
temporal, de la prueba técnica del prototipo completo, todavía limitada a
escenarios curados y ejecución controlada.

La Tabla~\ref{tab:9-e2e-capa1-capa2} resume la prueba reproducible posterior a la
integración de Capa 1. En ella, las latencias medias son 4.58844 ms para Capa 1,
52.16898 ms para Capa 2 y 0.30748 ms para Capa 3 degradada. Las métricas micro de
señales de Capa 1 son precisión 0.733333, sensibilidad 1.000000 y F1 0.846154. Estas cifras
proceden de cinco escenarios curados y solo acreditan coherencia funcional del flujo.

\begin{table}[ht]
\centering
\caption{Prueba reproducible Capa 1--Capa 2--Capa 3}
\label{tab:9-e2e-capa1-capa2}
\small
\begin{tabularx}{\textwidth}{|p{4.0cm}|p{1.8cm}|p{0.9cm}|X|}
\hline
\textbf{Caso} & \textbf{Prioridad} & \textbf{OK} & \textbf{Señales clave} \\ \hline
Incendio urbano con atrapados & P1 & Sí & atrapado, incendio, varias llamadas, riesgo vital \\ \hline
Accidente de tráfico grave & P1 & Sí & herido grave, atrapado, tráfico \\ \hline
Fuga química en nave industrial & P2 & Sí & intoxicación, riesgo vital textual \\ \hline
Inundación en garajes & P2 & Sí & meteo/inundación \\ \hline
Árbol caído en carretera secundaria & P3 & Sí & tráfico, sin falso atrapamiento \\ \hline
\end{tabularx}
\end{table}

\section{Evaluación exploratoria con participantes}

Como complemento de la evaluación técnica, se realizó un estudio exploratorio breve
para analizar la utilidad percibida y la comprensibilidad del prototipo. Se aplicó un
muestreo no probabilístico e intencional de nueve participantes, con acceso controlado
mediante Google Colab, casos representativos y un cuestionario de respuesta breve. El
instrumento empleó una escala Likert de 1 a 5 para valorar la razonabilidad de la
prioridad, la claridad de la explicación, la confianza aportada por las señales y la
visibilidad de la decisión humana, además de una pregunta abierta.

La muestra se distribuye en tres grupos de tres perfiles. El Grupo A reúne profesionales
vinculados a puestos de mando, comunicaciones críticas y emergencias; el Grupo B,
usuarios técnicos familiarizados con herramientas de IA; y el Grupo C, usuarios no
especializados en IA con distintos niveles de familiaridad operativa. Esta segmentación
permite contrastar perspectivas operativas, técnicas y de comprensibilidad general. El
Grupo A no se presenta como una muestra de operadores del 112: dos de sus perfiles se
especializan principalmente en comunicaciones y sistemas CIS, mientras que uno aporta
experiencia más próxima al despliegue de puestos de mando en emergencias. Los perfiles
funcionales anonimizados A1--C3 se detallan en el Anexo~\ref{cap:anexod}.

La Tabla~\ref{tab:9-evaluacion-exploratoria} presenta las medias observadas. La dimensión
mejor valorada es la claridad de que la decisión final corresponde a una persona, con
una media global de 4,63. La razonabilidad de la prioridad alcanza 4,14,
mientras que la claridad de la explicación y la confianza aportada por las señales
obtienen 4,04. Estos valores indican una percepción favorable del principio de
supervisión humana y una aceptación preliminar de la recomendación y su explicación.

\begin{table}[htbp]
\centering
\caption{Resultados agregados de la evaluación exploratoria con participantes}
\label{tab:9-evaluacion-exploratoria}
\small
\begin{tabularx}{\textwidth}{|X|c|c|c|c|}
\hline
\textbf{Dimensión evaluada} & \textbf{Grupo A} & \textbf{Grupo B} & \textbf{Grupo C} & \textbf{Global} \\ \hline
Prioridad recomendada razonable & 4,47 & 4,13 & 3,83 & 4,14 \\ \hline
Explicación clara & 4,17 & 4,37 & 3,60 & 4,04 \\ \hline
Señales o reglas aumentan la confianza & 4,50 & 4,23 & 3,40 & 4,04 \\ \hline
Decisión humana claramente identificada & 4,93 & 4,43 & 4,53 & 4,63 \\ \hline
\end{tabularx}
\end{table}

El análisis por perfiles aporta matices relevantes. El Grupo A valora especialmente la
supervisión humana (4,93) y la trazabilidad proporcionada por las señales
(4,50); el Grupo B otorga su mayor puntuación a la claridad de la explicación
(4,37); y el Grupo C presenta las valoraciones más bajas en señales o reglas
(3,40) y explicación (3,60). Los comentarios abiertos son coherentes con
este patrón: los perfiles operativos destacan la trazabilidad y la visibilidad de las
señales críticas, los perfiles técnicos solicitan una separación clara entre
priorización y explicación, y los no especializados demandan un lenguaje más directo.

La evaluación constituye evidencia descriptiva de utilidad percibida, no una validación
de aceptación ni una prueba operativa. El tamaño de la muestra, su selección intencional
y la ausencia de operadores del 112 impiden generalizar los resultados o inferir
efectos sobre el desempeño en una sala de emergencias.

\section{Limitaciones}

La evaluación presenta ocho limitaciones principales:

\begin{enumerate}
    \item La referencia P1--P4 es académica y débil, no una etiqueta operativa oficial
    del 112.
    \item El acuerdo de la supervisión débil es alto, pero las cuatro funciones comparten
    una puntuación heurística y no equivalen a anotadores independientes; tampoco sustituyen
    a una validación externa con personal experto del servicio.
    \item RuleFit-lite se ajusta con 3.000 de los 6.512 casos disponibles y la alternativa
    RuleFit externa con 50. Esta diferencia obliga a tratar \texttt{imodels} como prueba
    diagnóstica y no permite comparar eficiencia en condiciones equivalentes.
    \item El error de calibración esperado (\textit{Expected Calibration Error}, ECE)
    con objetivo \(ECE \leq 0{,}10\) queda pendiente de una evaluación específica
    \cite{Guo2017Calibration}.
    \item Las variables externas diferidas (AEMET histórico, SNCZI y Seveso)
    no participan en esta evaluación.
    \item Las pruebas integradas aportan evidencia de funcionamiento técnico P1--P4 y
    compatibilidad Capa 1--Capa 2 sobre 5 escenarios, pero su tamaño sigue siendo pequeño y la
    latencia media con LLM local, \(4{,}90\) s, requiere optimización antes de una
    demostración exigente.
    \item Los 7 errores P1--P4 de la partición de prueba estratificada son el principal foco de riesgo
    del modelo y deben analizarse individualmente antes de cualquier validación
    operativa externa.
    \item La evaluación exploratoria incluye nueve participantes seleccionados de forma
    intencional y no incorpora operadores del 112. Sus resultados describen percepciones
    de la muestra, pero no son estadísticamente generalizables ni acreditan aceptación
    operativa.
\end{enumerate}

Pese a estas limitaciones, la evaluación permite concluir que la Capa 2 seleccionada
obtiene una macro-F1 superior a la línea base heurística en las particiones evaluadas y mantiene la sensibilidad P1 por encima
del umbral fijado y respeta las restricciones de interpretabilidad y anti-leakage. Por
ello, RuleFit-lite se adopta como núcleo interpretable de Capa 2, siempre dentro del alcance académico del TFM. La evaluación se realiza
frente a etiquetas débiles P1--P4 y no frente a prioridades oficiales del 112, por lo
que la validación externa con personal experto sigue siendo necesaria antes de cualquier
uso operativo.

\clearpage
\begin{table}[p]
\centering
\caption{Síntesis final de la evaluación experimental}
\label{tab:9-sintesis-final}
\small
\begin{tabularx}{\textwidth}{|>{\RaggedRight\arraybackslash}p{3.0cm}|>{\RaggedRight\arraybackslash}p{4.2cm}|>{\RaggedRight\arraybackslash}X|}
\hline
\textbf{Elemento} & \textbf{Resultado principal} & \textbf{Interpretación y alcance} \\ \hline
Línea base heurística & 19 reglas; exactitud 0,694425; macro-F1 0,464284; sensibilidad P1 0,998473 & Referencia transparente y conservadora; detecta P1 con alta sensibilidad, pero discrimina peor entre las cuatro clases. \\ \hline
RuleFit-lite & Ajuste con 3.000 de 6.512 registros; exactitud 0,903648; macro-F1 0,905604; sensibilidad P1 0,909924; 30 reglas & Modelo seleccionado por su rendimiento de validación e interpretabilidad. Las métricas corresponden al estimador sin Puerta de Seguridad. \\ \hline
RuleFit externo & Ajuste diagnóstico con 50 registros; exactitud 0,768385; macro-F1 0,568239; sensibilidad P1 0,902778 & Contraste externo reducido. No permite inferir superioridad estadística ni computacional en condiciones equivalentes. \\ \hline
Prueba temporal 2022 & 735 casos; exactitud 0,884354; macro-F1 0,877594; sensibilidad P1 0,928767 & Primera evidencia de estabilidad temporal; no sustituye una validación prospectiva. \\ \hline
Falsos negativos P1 & 59 casos: 34 P1--P2, 18 P1--P3 y 7 P1--P4 & Los siete errores P1--P4 constituyen el principal riesgo y requieren especial atención. \\ \hline
Puerta de Seguridad & Posprocesamiento determinista comprobado en escenarios funcionales integrados & Salvaguarda separada de la evaluación estadística; su efecto no está incluido en las métricas offline de RuleFit-lite. \\ \hline
Evaluación exploratoria & Nueve participantes; medias globales entre 4,04 y 4,63 sobre 5 & Indicios descriptivos favorables, con mayor dificultad de comprensión de señales y reglas entre perfiles no especializados. \\ \hline
Limitaciones principales & Etiquetas académicas; funciones heurísticas correlacionadas; submuestreo de ajuste; contraste externo reducido; revisión interna conjunta & Los resultados son académicos y reproducibles, pero no acreditan validación operativa, acuerdo entre revisores independientes ni uso en un servicio 112 real. \\ \hline
\end{tabularx}
\end{table}
\clearpage
